\section*{Введение}
\addcontentsline{toc}{section}{Введение}

% Обосновывается актуальность разработки гибридной интеллектуальной системы
% для выбранной предметной области. Следует показать практическую или
% исследовательскую проблему.
%
% Формулируется цель курсовой работы: разработка гибридной интеллектуальной
% системы для решения определённого класса задач в выбранной предметной области.
%
% Перечисляются задачи, выполнение которых необходимо для достижения цели:
% - анализ предметной области, класса задач и существующих решений;
% - анализ базы знаний, решателя задач и интерфейса, разработанных
%   на предыдущих этапах курсового проектирования;
% - формирование требований к гибридной интеллектуальной системе;
% - проектирование пользователей, сценариев использования и архитектуры;
% - реализация, тестирование и демонстрация гибридной интеллектуальной системы.
%
% Кратко характеризуются применяемые подходы:
% - логико-символьный подход: база знаний, продукционные правила,
%   логический вывод, онтологии, семантические сети, поиск;
% - нейробионический подход: модели машинного обучения, нейронные сети,
%   языковые модели, эмбеддинги, классификация, ранжирование или генерация.
%
% Указывается практическая значимость системы, ожидаемый эффект применения,
% объект и предмет исследования, а также краткое содержание пояснительной записки.
